\section{Background}
\label{sec:prelim}

We define for the rest of the paper: learning to
reject for a fixed regression rule (Section~\ref{sec:bg-reject}) and the
standard PINN formulation for PDEs (Section~\ref{sec:bg-pinn}).

\subsection{Learning to reject in regression}
\label{sec:bg-reject}

Let $(X,Y)$ be a random pair with joint distribution $\Dcal$ on
$\Xcal \times \Ycal$, where $\Xcal \subseteq \R^{d_x}$ is a measurable
input space and $\Ycal \subseteq \R^{d_u}$ is a measurable response space.
Throughout,
$\psi:\Ycal\times\Ycal\to\R_{\geq 0}$ denotes a measurable non-negative
discrepancy, such as the squared loss
$\psi(a,b)=\tfrac{1}{2}\|a-b\|_2^2$; boundedness of $\psi$ is assumed only
when it is required by a stated consistency bound.

A regression rule with rejection is a pair $(h,r)$ in which
$h:\Xcal\to\Ycal$ is a measurable predictor and
$r:\Xcal\to\{0,1\}$ is a measurable rejector. We use the convention
$r(x)=0$ to accept the prediction $h(x)$ and $r(x)=1$ to abstain. Following
the classical rejection literature, abstention trades prediction error
against a rejection cost \citep{chow1970,bartlett2008classification,
cortes,cortes2024theory}. Most standard formulations take this cost to be
constant, i.e.\ $c(x)\equiv c$ \citep{chow1970,
bartlett2008classification,Geifman_El-Yaniv_2017,ni2019calibration,
cao2022generalizing,Mao_Mohri_Zhong_2023}; we write
$c:\Xcal\to[0,C]$ to allow the more general input-dependent case, expressed
in the same units as $\psi$. We work in a second-stage
predictor--rejector setting \citep{cortes,Mao_Mohri_Zhong_2023,
cheng2023regression}, where $h$ is fixed and only $r$ is learned. The
pointwise loss is
\begin{equation}
  \ell_{\mathrm{abs}}\bigl(h,r;\,x,y\bigr)
  :=
  \psi\bigl(h(x),y\bigr)\ind\{r(x)=0\}
  +
  c(x)\ind\{r(x)=1\},
  \label{eq:abs-loss}
\end{equation}
with population risk
$\Ecal_{\ell_{\mathrm{abs}}}(h,r)
:=\E_{(X,Y)\sim\Dcal}[\ell_{\mathrm{abs}}(h,r;\,X,Y)]$. When the rule
accepts, it pays the discrepancy of the fixed predictor; when it abstains,
it pays $c(X)$.

\begin{lemma}[Bayes-optimal rejector]
\label{lem:bayes-rej}
For a fixed predictor $h$, define
$\eta_h(x):=\E[\psi(h(x),Y)\mid X=x]$. Then, for $\Pr_X$-a.e.\ $x$, the
Bayes-optimal rejector is
\[
  r^\star(x) \in
  \begin{cases}
    \{0\}, & \eta_h(x) \leq c(x),\\
    \{1\}, & \eta_h(x) > c(x),
  \end{cases}
\]
where $r^\star(x)=0$ accepts $h(x)$ and $r^\star(x)=1$ abstains.
\end{lemma}

The proof is the pointwise comparison of the two available costs,
$\eta_h(x)$ and $c(x)$. Thus the rejector only needs to learn the sign of
$\eta_h(x)-c(x)$. We use a measurable score $s:\Xcal\to\R$ from a rejector
class $\Hcal_r$ and define $r_s(x):=\ind\{s(x)>0\}$. A smooth surrogate
$\Phi_{\mathrm{abs}}$ for the discontinuous loss
\eqref{eq:abs-loss} is useful only if small surrogate excess risk controls
the true abstention excess risk in $\Hcal_r$, in the sense of
$\Hcal$-consistency bounds
\citep{Awasthi_Mao_Mohri_Zhong_2022_multi,
mao2023crossentropylossfunctionstheoretical,Mao_Mohri_Zhong_2023}. This
abstract setup is the only learning-to-reject machinery needed in the sequel.

\subsection{Physics-informed neural networks}
\label{sec:bg-pinn}

Let $\Omega\subset\R^d$ be a bounded open domain with Lipschitz boundary
$\partial\Omega$, and consider the boundary-value problem
\begin{equation}
  \Ncal[u](x)=f(x), \quad x\in\Omega,
  \qquad
  \mathcal{B}[u](x)=g(x), \quad x\in\partial\Omega,
  \label{eq:pde}
\end{equation}
where $\Ncal$ is a possibly nonlinear differential operator,
$\mathcal{B}$ is a boundary operator, and $f,g$ are given data in function
spaces appropriate for the problem. For the strong-form notation used here,
we assume \eqref{eq:pde} admits a unique sufficiently smooth solution
$u^\star:\overline{\Omega}\to\R^{d_u}$, with enough regularity for the
operators below to be evaluated pointwise.

A PINN is a parameterised map
$\widehat{u}_\theta:\Omega\to\R^{d_u}$, typically a feed-forward neural
network with differentiable activations, trained so that the PDE and
boundary residuals are small. The residuals are
\begin{equation}
  \mathfrak{r}_{\mathrm{int}}^\theta(x)
  :=
  \Ncal[\widehat{u}_\theta](x)-f(x),
  \qquad
  \mathfrak{r}_{\mathrm{bd}}^\theta(x)
  :=
  \mathcal{B}[\widehat{u}_\theta](x)-g(x),
  \label{eq:residuals}
\end{equation}
and are computed by automatic differentiation through
$\widehat{u}_\theta$ when the operators are represented in strong form
\citep{raissi2019physics,karniadakis2021physics,lu2021deepxde}. Given
interior collocation points
$\{x_i^{\mathrm{int}}\}_{i=1}^{N_{\mathrm{int}}}\subset\Omega$ and boundary
collocation points
$\{x_j^{\mathrm{bd}}\}_{j=1}^{N_{\mathrm{bd}}}\subset\partial\Omega$, a
standard PINN objective is
\begin{equation}
  \Lcal_{\mathrm{PINN}}(\theta)
  :=
  \frac{\lambda_{\mathrm{int}}}{N_{\mathrm{int}}}
  \sum_{i=1}^{N_{\mathrm{int}}}
    \bigl\|\mathfrak{r}_{\mathrm{int}}^\theta(x_i^{\mathrm{int}})\bigr\|^2
  +
  \frac{\lambda_{\mathrm{bd}}}{N_{\mathrm{bd}}}
  \sum_{j=1}^{N_{\mathrm{bd}}}
    \bigl\|\mathfrak{r}_{\mathrm{bd}}^\theta(x_j^{\mathrm{bd}})\bigr\|^2,
  \label{eq:pinn-loss}
\end{equation}
with non-negative weights $\lambda_{\mathrm{int}}$ and
$\lambda_{\mathrm{bd}}$. Data-fit terms on observed solution values can be
added analogously.

A central difficulty is that \eqref{eq:pinn-loss} directly constrains
residuals only through sampled empirical averages, whereas the deployment
quantity of interest is often the local solution error
$\widehat{u}_\theta(x)-u^\star(x)$. A-posteriori results connect these
quantities at the level of global norms. Schematically, under suitable
stability, regularity, approximation, and sampling assumptions,
\citet{mishra2022estimates} obtain bounds of the form
\begin{equation}
  \bigl\|\widehat{u}_\theta-u^\star\bigr\|_{L^2(\Omega)}
  \leq
  C_\Ncal
  \bigl\|\mathfrak{r}_{\mathrm{int}}^\theta\bigr\|_{L^2(\Omega)}
  +
  C_\mathcal{B}
  \bigl\|\mathfrak{r}_{\mathrm{bd}}^\theta\bigr\|_{L^2(\partial\Omega)}
  +
  \mathrm{approximation/sampling\ terms},
  \label{eq:apost}
\end{equation}
where the constants and extra terms depend on the PDE class, the domain,
and the sampling scheme. Equation~\eqref{eq:apost} should therefore be read
as a global stability statement, not as a pointwise certificate of trust.
The abstention layer introduced in Section~\ref{sec:method} addresses this
local question by taking the fixed predictor $h$ in
Section~\ref{sec:bg-reject} to be the frozen PINN
$\widehat{u}_\theta$, whether it was trained by \eqref{eq:pinn-loss} or by
modern variants such as adaptive weighting, causal training, adaptive
sampling, or domain decomposition \citep{wang2021understanding,
mcclenny2023self,wang2022respecting,wu2023comprehensive,daw2023mitigating,
jagtap2020extended,jagtap2020conservative,moseley2023finite}.
